\documentclass[../main.tex]{subfiles}
\begin{document}
%==========================================TIÊU ĐỀ TẬP========================================
\begin{table}[h]
  \begin{adjustwidth}{-1cm}{}
    \begin{tabular}{>{\centering\arraybackslash}p{6cm}|>{\raggedright\arraybackslash}p{12.5cm}}
      \multirow{5}{6cm}
      % Cột bên trái: ÉPISODE và số tập
      {\\[-40pt]
        \flaregothic\fontsize{20pt}{20pt}\selectfont \centering
        {\contour{errorfour}{\textcolor{black}{ÉPISODE}}}\color{black}\\
        \burbank\fontsize{72pt}{72pt}\selectfont
        \includegraphics[height=3.12359cm]{../../../../Image/Book?/number/num20.png}
        \\[18pt]
        % \flaregothic\fontsize{14pt}{14pt}\selectfont
        % \color{epattr}BIENVENUE, \uline{\color{Banpire} Mel} !
        % \flaregothic\fontsize{18pt}{18pt}\selectfont
        % \color{epattr}ÉPISODE PERDU
        \color{black}
      }
       &
      % Dòng thứ nhất: tên tập tiếng Nhật
      {\fotskip\fontsize{18pt}{18pt}\selectfont \ruby{幸}{しあわ}せな\ruby{世}{せ}\ruby{界}{かい}}
      \\[6pt]
      % Dòng thứ hai: tên tập tiếng Anh
       & {\cafeteria\fontsize{18pt}{18pt}\selectfont A Happy World}                              \\[4pt]
      % Dòng thứ ba: tên tập tiếng Việt
       & {\vnmsans\fontsize{18pt}{18pt}\selectfont Thế giới hạnh phúc giả tạo}                   \\[8pt]
      % Dòng thứ tư: Nhân vật xuất hiện
       & {\caviar{\fontsize{14pt}{14pt}\selectfont Track used: The Samsara (Sewind - FindersX)}}
      \\[6pt]
      % Dòng cuối cùng: Thời gian diễn ra của tập (thường sẽ là ngày phát hành)
       & {\continuum\fontsize{16pt}{16pt}\selectfont Ngày phát hành: 02/07/2022}                 \\[8pt]
    \end{tabular}
  \end{adjustwidth}
\end{table}
%===========================================NỘI DUNG==========================================
\color{black}

\pov{Hikawa Kuon}

\dots

\dots

``\dots Ưm\dots''

Buổi sáng hôm đó, ánh nắng xuyên qua lớp rèm cửa, chạm vào gò má tôi bằng một thứ ấm áp, nhưng không đủ nóng để làm da thịt thực sự cảm thấy sống. Mọi thứ xung quanh đều quá yên bình đến mức bất thường.

Tôi ngồi dậy, nhìn quanh căn phòng nhỏ: bàn học gọn gàng, giá sách ngăn nắp đến mức như vừa được ai đó sắp đặt. Không có lấy một hạt bụi, không có dấu hiệu nào cho thấy nơi này từng có người ở thật. Cánh cửa tủ quần áo khép hờ, nhưng không hề có mùi vải cũ, cũng chẳng có vết nhăn nào trên chiếc áo sơ mi trắng được treo thẳng tắp.

Tôi đưa tay vuốt qua mặt bàn: lạnh, trơn, như thể vừa được lau bằng một chiếc khăn ẩm hoàn hảo ngay trước khi tôi mở mắt.

``\dots Đây là\dots\ đâu đây?''

Tiếng thì thầm của tôi vang lên, rồi chìm vào bức tường phủ sơn trắng tinh, không dấu vết. Tôi bước ra ngoài ban công. Con đường dẫn đến trường rợp bóng cây, những chiếc lá đung đưa trong làn gió nhẹ, và những người học sinh mỉm cười, chào hỏi nhau như trong một bức tranh tĩnh vật. Không một tiếng còi xe, không ai vội vã, thậm chí cả tiếng gió cũng đều đặn một cách máy móc, một nhịp hoàn hảo không trật lấy một ly.

``Thế giới hoàn hảo, hử\dots''

Tôi lẩm bẩm, khẽ siết tay nhìn vào chiếc đồng hồ nơi Game trú ngụ. Kim giây vẫn chạy, đều đặn và vô cảm. Tôi có cảm giác như thời gian ở đây chỉ là một phần trong trò diễn, giống như mọi thứ được lập trình sẵn để đánh lừa cảm giác của con người.

Tôi thử nhắm mắt, đếm ba nhịp thở; khi mở mắt ra, đám mây trên trời đã dịch chuyển đúng hai ly về phía đông, như thể ai đó đã căn chỉnh bằng thước thủy. Trong lòng tôi bỗng dâng lên một cảm giác trống rỗng kỳ lạ: không phải nỗi sợ, cũng không phải buồn bã, mà là thứ cảm giác khi bạn biết mình đang đứng giữa một bức tranh 3D khổng lồ, và chỉ cần chạm vào là cả thế giới sẽ vỡ vụn thành từng mảnh nhựa mỏng.

``Game-san,'' tôi gọi khẽ, ``ông có cảm thấy chỗ này\dots\ quá bất thường không?''

Giọng ông vang lên, vọng từ chiếc đồng hồ như một tiếng vọng xa xăm. ``Yên tĩnh đến mức giả tạo. Cẩn thận đấy, Kuon-user. Đây không phải thế giới thực. Mọi thứ ở đây chỉ đang diễn lại một kịch bản thực thể `Sorra' đó muốn cậu tin.''

Tôi khẽ cười, ánh nhìn dừng lại ở những tòa nhà xa xa đang phản chiếu ánh nắng như gương. Tôi biết cái ``kịch bản'' ông ta nói là gì: một chuỗi sự kiện lặp đi lặp lại, chìm trong thứ gọi là hạnh phúc trong mộng tưởng mà tôi, Kaigou và Suzuki đã từng trải qua ở những vòng lặp trước. Nhưng lần này, dù đã cảnh giác, tôi vẫn cảm thấy một luồng hơi ấm nhẹ chạy dọc sống lưng.

Tôi siết chặt tay vịn ban công, cảm giác sắt sơn lạnh buốt len vào lòng bàn tay, nhưng tôi không buông. Dù biết rõ là ảo ảnh, tim tôi vẫn đập nhanh hơn, giống như linh cảm rằng, đâu đó trong cái thế giới hoàn hảo này, Shino đang tồn tại, không phải như một mảnh ghép được gắn vào kịch bản, mà như một điểm nóng duy nhất còn sót lại, thứ mà ngay cả kẻ viết kịch bản cũng không thể xóa nhòa.

\pov{Hikawa Kaigou}

``Suzu,'' tôi gọi khẽ, giữa tiếng dầu xèo xèo, ``em có cảm thấy\dots\ mùi này hơi giống\dots\ mùi ở căn bếp cũ không?''

Em ấy khựng đũa một giây, rồi nhẹ nhàng lật miếng trứng. ``Em không nhớ nữa, Onee-sama. Có lẽ chỉ là\dots\ mùi bình thường thôi ạ.''

Tôi nhìn xuống chảo trứng mà em ấy đang rán. Lòng trắng tròn vo, lòng đỏ căng mọng, không một vết rách. Giống như ai đó đã đo lường kỹ lưỡng nhiệt độ, thời gian, và cả độ rung tay.

``Ừ, bình thường,'' tôi đáp, nhưng trong đầu lại vang lên tiếng còi xe buýt năm cấp hai, mùi khét của bánh mì cháy vào một buổi sáng ăn vội.

Tôi dậy sớm hơn thường lệ. Không rõ vì sao, nhưng cơ thể lại phản ứng theo thói quen cũ, như thể đã từng sống ở nơi này từ rất lâu. Bên ngoài cửa sổ, bầu trời loang màu cam nhạt của bình minh, thứ ánh sáng mà tôi từng nghĩ là dấu hiệu của một ngày mới. Nhưng hôm nay, nó chỉ khiến tôi thấy rờn rợn.

Suzu đã có mặt trong bếp, chiếc tạp dề trắng phủ lên đồng phục, tay đảo nhẹ chảo trứng trên bếp. Em ấy trông điềm đạm như mọi khi, nhưng ánh mắt thì có chút xa xăm.

Tôi kéo cánh cửa bếp, bước xuống với chiếc cà vạt buông lơi, ánh mắt hiện rõ sự cảnh giác cao độ. Tôi hiểu cảm giác của cậu ấy là gì: sự yên bình đến mức đáng ngờ đó khiến đến chúng tôi cũng không tài nào tin nổi.

``Cậu dậy muộn quá đó,'' tôi buông lời trêu, cố phá tan bầu không khí căng thẳng.

Em tôi quay lại, đôi môi khẽ cong, giọng nhẹ mà đều: ``Onii-sama, nhanh lên để còn tới trường nào. Muộn học rồi đấy ạ.''

Tôi bước lại gần, đặt mấy lát bánh mì lên đĩa, cảm nhận hơi nóng của lò nướng vẫn còn âm ấm. Mọi thứ ở đây đều quá hoàn hảo: căn bếp sáng bóng, hương bơ thoang thoảng trong không khí, đồng hồ treo tường nhích từng nhịp đều đặn. Từng chi tiết nhỏ dường như đều được sắp đặt để trấn an chúng tôi rằng đây là đời thật.

Suzu đặt đĩa trứng xuống trước mặt tôi, mỉm cười. ``Mời cả nhà ăn cơm,'' cả ba chúng tôi chắp tay đồng thanh.

Tôi cầm dĩa lên, nhưng không dám cắt. Sợ rằng nếu mảnh trứng vỡ ra, cả thế giới này cũng sẽ nứt theo.

``Này, Kuon-kun,'' tôi khẽ gọi, cố tình dùng giọng hồn nhiên nhất có thể, ``hôm nay cậu thấy tay nghề làm bếp của Suzu thế nào? Trứng rán hoàn hảo thế này, chắc phải có bằng cấp đấy nhỉ?''

Kuon-kun ngẩng lên, khóe môi khẽ nhếch, một chút ấm áp hiếm hoi lướt qua đôi mắt cậu. ``Ừ, ngon thật. Mà Suzu, em nhớ cho thêm chút tiêu nhé, nó sẽ thơm hơn đấy.''

Suzu cười khúc khích, đôi má ửng hồng. ``Vâng, em nhớ rồi. Lần này em cho thêm tiêu xanh nữa, Onii-sama thích mà.''

Trong phút chốc, căn bếp như ngập tràn hơi ấm thật sự. Tôi cảm thấy tim mình đập chậm lại, như thể mình đang sống trong một khoảnh khắc không bị lập trình. Nhưng rồi, tiếng tích tắc của đồng hồ treo tường lại gõ vào tai tôi, đều đặn, vô cảm, nhắc nhở rằng khoảnh khắc này cũng chỉ là một phần của kịch bản.

Tôi rót sữa ra cốc, nhìn làn sương sớm ngoài cửa sổ. Dòng suy nghĩ lặng lẽ trôi qua trong đầu: ``Nếu đây cũng chỉ là một bản sao khác của thế giới đó\dots\ thì lần này, liệu Shino có thể nhận ra?''

``Kuon-user,'' giọng Game vang lên từ chiếc đồng hồ, trầm và đều, ``tôi hiểu cảm giác của cậu. Nhưng trước mắt thì, mọi người cứ sống như bình thường đã. Đừng để kẻ viết kịch bản thấy chúng ta đã nhận ra.''

Kuon-kun siết nhẹ cốc sữa, cảm thấy hơi nóng lan tỏa qua lòng bàn tay. ``Ý ông là\dots\ chúng ta nên giả vờ bị lừa?''

``Không phải giả vờ,'' Game đáp, ``mà là chọn thời điểm. Khi mọi thứ đang diễn theo đúng nhịp, bất kỳ sai lệch nào cũng có thể khiến cỗ máy tái khởi động. Hãy để họ nghĩ rằng chúng ta vẫn đang mơ.''

Suzu ngẩng lên, ánh mắt thoáng lo âu. ``Vậy em\dots\ vẫn cứ rán trứng như thế này?''

``Ừ,'' tôi mỉm cười an ủi, ``và nhớ cho thêm tiêu nhé. Chúng ta cần giữ lấy những gì có thể giữ, kể cả khi biết nó không thật.''

Tôi gật đầu, rồi cắt một miếng trứng nhỏ, đưa lên miệng. ``Vậy có lẽ\dots\ chúng ta cứ sinh hoạt như mọi khi ta hay làm thôi nhỉ?''

Tôi nhìn hai người họ, cảm thấy một thứ gắn bó mới, không phải của máu mủ, mà của những kẻ cùng bị mắc kẹt trong một vở kịch. Và rồi, tôi cũng cắt miếng trứng, để nó vỡ ra thành những mảnh vàng ươm, tan trong miệng, như một lời hứa rằng: dù thế giới này có là giả, tình cảm chúng tôi dành cho nhau vẫn là thật.

``Rồi sẽ tới cảnh buổi sáng yên bình,'' tôi lẩm bẩm, ``rồi đến trường, gặp vài người bạn mới, nghe tin về `trận cúm' ở ga tàu. Chiều về nhà, rồi sáng hôm sau lặp lại y hệt.''

Kuon-kun nghe thế, chỉ im lặng. Cậu đặt chảo xuống, lau tay bằng khăn, rồi nói nhỏ như một tiếng thở dài: ``Và kết thúc bằng một vụ động đất.''

Tôi ngẩng lên nhìn cậu, bắt gặp ánh mắt hiểu rõ mọi chuyện hơn cả lời nói.

Chúng tôi đều nhớ vòng lặp ấy, lần đầu tiên khi chúng tôi bị mắc kẹt ở đây. Khi mà mọi thứ sụp đổ, rồi tái tạo lại, chỉ để tiếp tục lặp đi lặp lại một giấc mơ hạnh phúc giả tạo.

Tôi rót sữa ra cốc, nhìn làn sương sớm ngoài cửa sổ. Dòng suy nghĩ lặng lẽ trôi qua trong đầu: ``Nếu đây cũng chỉ là một bản sao khác của thế giới đó\dots\ thì lần này, liệu Shino có thể nhận ra?''

\pov{Hikawa Suzuki}

Ăn sáng xong, chúng tôi nhanh chóng chuẩn bị cho bữa trưa ở trường, rồi rời khỏi nhà khi cả thị trấn bắt đầu đón những tia nắng đầu tiên.

Mặt trời lên cao hơn một chút, phủ xuống thị trấn Aogami một lớp ánh sáng dịu màu vàng nhạt. Tôi bước chậm hơn một chút so với Onee-sama, trong khi Onii-sama đi trước. Chúng tôi đang trên đường tới trường, men theo con dốc quen thuộc, nơi gió núi thổi mang theo mùi gỗ thông và hương cỏ mới cắt.

Thị trấn này\dots\ thật đẹp.

Aogami nằm giữa những dãy núi xanh thẳm, những ngôi nhà gỗ cổ vẫn đứng nép bên cạnh những tòa bê tông mới xây. Con đường lát đá uốn quanh sườn dốc, dẫn đến một dãy cửa hàng nhỏ: tiệm bánh, quán cà phê, hiệu sách. Dưới chân núi, chiếc xe điện sơn màu bạc chạy êm ru, phản chiếu ánh sáng như đang trôi giữa một thế giới không tì vết.

Hồi tôi còn là ``Reiko'', tôi từng yêu khung cảnh này - một Aogami dung dị và ấm áp. Nhưng giờ, khi nhìn lại, tôi chỉ thấy nó hoàn hảo đến mức giả tạo. Không một chiếc lá rơi lệch hướng, không một âm thanh thừa thãi. Mọi người đều mỉm cười, chào hỏi nhau như thể chẳng có gì tồi tệ từng xảy ra. Một sân khấu được dựng nên đầy giả tạo, tới mức độ những người bình thường như chúng tôi còn phải thấy nhướng mày.

Kuon khẽ nói, giọng anh trầm hơn bình thường: ``Mọi thứ ở đây đều yên bình tói mức kỳ lạ.''

Kaigou gật đầu, mái tóc chị tôi khẽ lay trong gió. ``Một Aogami hoàn hảo hơn cả trí nhớ. Nhưng cũng chính vì thế mà nó không hề thật một chút nào.''

Tôi không đáp, chỉ nhìn những dãy cột điện nối nhau qua các mái nhà, những sợi dây chằng chịt như mạng lưới trói buộc ký ức. Trong đầu tôi, hình ảnh Shino-san thoáng hiện. Cô ấy giờ đang ở đâu? Và liệu cô có nhận ra đây chỉ là một vòng lặp khác, một ảo ảnh được vẽ ra để nuốt chửng nỗi đau?

Tôi mím môi, nắm chặt quai cặp. ``\textit{Nếu đây là thế giới của Sora,}'' tôi nghĩ thầm, ``\textit{thì Shino-san sẽ phải đối mặt với chính mình. Nhưng liệu cô ấy có đủ can đảm để làm điều đó không?}''

Ba chúng tôi tiếp tục bước đi, bóng đổ dài trên con đường lát đá sáng bóng. Mọi thứ vẫn yên bình đến mức kỳ dị, cho đến khi\dots\ tôi thấy dáng người quen thuộc ở phía xa.

Một cô gái tóc nâu với chiếc cài tóc hình ngôi sao với lông đỏ lệch bên trái đang thong dong đi về phía trước, những bước chân không có vẻ gì như đang bị gượng ép cả.

Misora Shino-san. Chính xác hơn, là bản thể ``Sora.''

Cô quay lại, nụ cười nhẹ hệt như chưa từng có gì xảy ra, một nụ cười khiến tim tôi lạnh buốt.

Tôi nghe Onee-sama thì thầm, giọng cô khàn đi: ``Chính là cô ấy.''

Và khi ánh mắt của cô ấy gặp lấy chúng tôi, điều đó có nghĩa rằng\dots\ vòng lặp này đã thật sự bắt đầu.

\pov{-}

``Chào buổi sáng, Kyuen-san, Saikyou-san, Reiko-san!''

Giọng Shino vang lên trong trẻo, nhẹ đến mức như hòa vào làn gió sớm của thị trấn. Cô mỉm cười tươi rói, vẫy tay chào ba người họ như thể đã thân quen từ lâu, một hành động tưởng chừng nhỏ, nhưng lại khiến cả ba người họ Hikawa khựng lại.

Ba cái tên ấy\dots

Không phải Kuon, Kaigou hay Suzuki. Mà là Kyuen, Saikyou, Reiko, những cái tên thuộc về họ trong vòng lặp đầu tiên, khi mọi thứ chỉ vừa bắt đầu, và họ vẫn còn ngây thơ tin rằng mình đang sống một cuộc đời bình thường.

Trong vài giây, không ai nói gì. Chỉ có gió thoảng qua, mang theo tiếng chim ríu rít trên mái nhà, hòa lẫn cùng nhịp bước chân của những người dân đang bắt đầu một ngày mới.

Rồi giọng của Game vang lên, không từ đâu cả, mà từ chiếc đồng hồ trên cổ tay Kuon: ``Tốt nhất là nên cư xử như bình thường. Cô ta sẽ không để tâm đâu.''

Chỉ một câu ngắn ngủi, nhưng đủ để Kuon hiểu. Anh liếc nhìn hai người còn lại, cả ba khẽ gật đầu, như một tín hiệu im lặng. Nếu Shino gọi họ bằng những cái tên đó, thì họ cũng sẽ nhập vai như thể đó là sự thật.

``Chào buổi sáng, Shino-san,'' Kuon đáp, giọng điềm đạm. ``Hôm nay trông cậu có vẻ tỉnh táo hơn đấy.''

Shino cười hồn nhiên. ``Fufu. Hôm qua tớ suýt muộn đó nha! Còn ngủ gật lúc cuối buổi nữa chứ\dots\ Nhưng hôm nay thì khác rồi, tớ dậy sớm hơn, có nhiều thời gian chuẩn bị kỹ hơn, và còn--''

``Vẫn liến thoắng như mọi khi nhỉ,'' Kaigou xen vào, cố giấu đi nụ cười nhạt. ``Ít nhất cũng phải để người khác nói chứ.''

``Onee-sama vẫn cứ thẳng tính như mọi khi nhỉ,'' Suzuki cũng phì cười.

``Saikyou-san thật là\~{}! Reiko-san, cậu ấy bắt nạt mình kìa\~{}!'' Shino chu môi, nhưng rồi cả nhóm cùng bật cười.

Con đường dẫn tới nhà ga hôm nay rực sáng bởi nắng sớm. Dưới ánh sáng ấy, thị trấn Aogami trông như một bức tranh hoàn hảo: từng hàng cây, từng viên đá lát đường, từng tòa nhà nhỏ xinh đều hiện ra sắc nét đến mức quá mức thật. Trong mắt Kuon và Kaigou, khung cảnh ấy không còn là niềm an yên; nó là dấu hiệu của một chu trình hoàn hảo, nơi mọi thứ được ``thiết kế'' để tái diễn vô tận.

Kuon khẽ liếc sang hàng cây phong rợp bóng bên vỉa hè: lá xanh non đứng im như tạc, không chuyển động dù gió đang lùa qua từng kẽ áo. Một giây sau, tất cả lá bỗng rung lên đồng loạt theo nhịp giật cố định - ba lần, đều như máy, rồi lại bất động. Không một chiếc lá nào rơi lệch quỹ đạo đã định sẵn.

Phía trước, một bà cụ bán bánh rán đang cúi xuống thu dọn chiếc ghế gỗ; bóng của bà trải dài trên nền đá. Bất ngờ, cái bóng nhòe thành vệt mờ đen, như ai đó xóa vội bằng khăn ướt, rồi trong nháy mắt tái hiện nguyên vẹn, góc nghiêng vẫn khớp đến từng chuyển động. Kaigou nắm chặt tay áo Kuon, ra hiệu ``đừng nhìn lâu''.

Từ phía trường học, tiếng chuông reo vang, thánh thót. Đồng hồ trên cổ tay Kuon chỉ 06:20, chưa đến giờ vào lớp. Tiếng chuông dứt bốn hồi, im bặt, rồi lặp lại y hệt 06:20:32, như bản phát lại bị kẹt. Kuon đếm thầm; đến lần thứ ba thì tiếng chuông mới chịu im, và mặt trời bất ngờ nhảy lên một khúc cao hơn, như thời gian bị chỉnh nhanh để bắt kịp lịch trình.

Những đường lỗi ấy thoắt ẩn thoắt hiện, chỉ đủ khiến bụng họ lạnh toát, rồi biến mất. Cửa hiệu bánh vẫn tỏa mùi vani ngọt sắc; chú mèo tam thể vẫn nằm cuộn tròn trên bậc thềm, đuôi khẽ đập đều, đúng bốn lần mỗi mươi giây, không thừa không thiếu. Dãy núi phía xa được phủ một lớp sương mỏng, đường viền sắc cạnh như được cắt bằng lưỡi dao lam. Không ai trong thị trấn - ngoại trừ ba người họ - dường như nhận ra những đứt gãy vi mô ấy của thực tại.

Shino bước thong dong, mái tóc nâu khẽ đung đưa theo nhịp bước hoàn toàn đúng chuẩn từng bước một. Mỗi lần cô quay lại cười, ánh mắt lấp lánh nhưng không phản chiếu, như thể phía sau con ngươi chỉ là một màn ánh sáng mờ đục. Họ tiếp tục đi, bước chân vẫn đều, vờ như không thấy, để giữ cho bức tranh Aogami không bị rách thêm bất kỳ một đường lỗi nào nữa.

\ct

Trước cửa nhà ga, năm cô gái đang đứng đợi.

Shino nhanh nhảu giơ tay chào, nụ cười cô rạng rỡ như nắng. ``Chào buổi sáng, mọi người!''

\img{e3-4a}

Ngay giây đó, ánh mắt của ba người mang họ Hỏa Xuyên giao nhau.

Năm người, lại là năm người.

Hệt như mọi vòng lặp trước: năm gương mặt, năm cái tên, năm hình bóng khác nhau nhưng cùng một khuôn mẫu.

Trong đầu Kuon, những ký ức lướt qua như những thước phim tua nhanh:

Inari, Hanabi, Saki, Touka, Suzune tại vòng lặp đầu tiên.

Nanase, Sakura, Yae, Shiina, Honoka từ thời quá khứ, nơi mà lời nguyền về trường trung học bắt đầu.

Akane, Yuka, Miku, Saya, Kaoru xoay quanh nhà ga.

Và bây giờ, là Mitsuki, Yuki, Hotaru, Kana, Kaande. Hay trong con mắt của một tổng quản lý, họ mang dáng dấp như Luna, Roboco, Flare, Towa và Kanata.

Những người bạn học khác nhau trong mỗi thế giới, nhưng luôn cùng giữ một cấu trúc, một hằng số bí ẩn không thể sai lệch.

Kaigou khẽ thì thầm: ``\vie{Lại là năm người\dots\ Có vẻ như khuôn mẫu không thay đổi.}''

``\vie{\dots Ừ,}'' Kuon đáp, mắt anh nhìn xa xăm. ``\vie{Mọi thứ bắt đầu lại, chỉ là bằng một lớp vỏ khác.}''

``\vie{Bản chất thì\dots\ vẫn không thay đổi ạ?}'' Suzuki nhỏ giọng, nhìn năm người bạn học đó. Cả hai anh chị cùng gật, không nói lấy một câu.

Chiếc đồng hồ trên cổ tay Kuon rung nhẹ. Giọng Game vang lên, trầm và chậm rãi như tiếng chuông ngân dưới đáy giếng: ``Rõ ràng nó có kịch bản được dựng sẵn. Năm bản sao, năm vai diễn, năm cái tên, nhưng chỉ có một mẫu số: `tình bạn đầu đời'. Cô ta không sáng tạo ra họ, cô ta chỉ sao chép cấu trúc cảm xúc mà chính cô ta khao khát.''

Kuon thấy lòng bàn tay mình ướt đẫm mồ hôi. Cậu nhớ lại cảm giác lần đầu tiên được gọi là ``Kyuen'' trong vòng lặp đầu, cậu cảm thấy có gì đó ngờ ngợ nhưng rồi cũng dần chấp nhận. Giờ đây, khi nghe lại ba âm tiết ấy, cậu chỉ thấy lạnh, như có ai đang tái hiện ký ức ngọt ngào rồi nhét một viên đá vào chỗ trống.

Kaigou siết chặt tay cầm cặp, đốt ngón tay trắng bệch. Cô nhìn thấy bản thân mình - không, là ``Saikyou'' của những ngày tháng đầu - đứng giữa hành lang trường, được bạn bè vây quanh, cười đến nỗi quên mất mình đang diễn. Cảm giác đó bây giờ trở thành một chiếc vòng cổ siết chặt: càng giật mạnh, càng nhớ rõ mình từng yêu thích thứ ảo ảnh ấy đến nhường nào.

Suzuki cúi đầu, mái tóc che khuất mắt. Cô thấy ``Reiko'' trong ký ức - cô bé với một sức hút tuyệt vời được các bạn học mếm mộ. Bây giờ, khi biết rằng khoảnh khắc ấy chỉ là một đoạn mã được gọi lại, cô không thể để bản thân mình bị lừa nữa.

Ba người họ cùng nhìn lên bầu trời. Mặt trời vẫn ở đó, hoàn hảo đến mức không một tia sáng nào lệch góc. Và họ hiểu: thế giới này không chỉ là một vở kịch, nó là bản giao hưởng do chính nỗi nhớ của kẻ đạo diễn viết nên, dùng những con người thay thế cho nốt nhạc. Họ không phải khán giả, cũng chẳng phải diễn viên; họ là những nốt nhạc bị kẹt giữa hai phách, mãi mãi không được phép dứt ra.

Trong lúc họ nói chuyện, Kana - cô gái với mái tóc tím xoắn tít buộc hai bên - quay lại, vẫy tay gọi to: ``Ê, mấy người đứng đó làm gì thế? Mau lên nào, sắp trễ rồi!''

Câu nói ấy kéo họ trở lại thực tại. ``Đ-Được rồi! Bọn tôi tới liền!''

Shino cười, chạy tới nhập nhóm cùng năm cô bạn kia, vừa cười vừa nói chuyện rôm rả, đuôi mắt cô khẽ liếc về phía cổng trường như thể đang đếm từng bậc thang để chắc rằng ước mơ về một cuộc sống học đường bình yên, mà cô không hề hay biết là đã từng sụp đổ bao nhiêu lần.

Còn Kuon, Kaigou và Suzuki bước theo sau, ánh mắt xen lẫn giữa sự dè chừng và bình thản, những kẻ đã nhìn thấy kẽ nứt sau lớp vỏ hoàn mỹ, và quyết định bước tiếp, chỉ để chờ thời khắc gảy lên nốt nhạc sai nhịp đầu tiên.

Một vòng lặp mới bắt đầu, ngày đầu tiên đến trường, một lần nữa.

Nhưng lần này, họ không còn là những học sinh ngây thơ.

Họ là những kẻ đang theo dõi thế giới từ bên trong.

\ct

\begin{center}
  \uline{Lớp 1-C - Trường Trung học Aogami\\Giờ nghỉ giữa giờ.}
\end{center}

Buổi sáng hôm ấy trôi qua với tốc độ nhanh đến mức cả bốn người vẫn chưa kịp định thần.

Từ tiết Ngữ văn đầu tiên với thầy giáo trung niên nghiêm nghị, khiến cả lớp im phăng phắc; đến tiết Toán thứ hai, nơi cô giáo phát một bài kiểm tra ngắn để thăm dò trình độ. Kaigou và Kuon, như thường lệ, vẫn hăng hái làm bài, thậm chí còn nhanh đến mức khiến vài bạn phải ngước nhìn.

Rồi tiết Lịch sử thứ ba với thầy giáo trẻ tuổi, giọng kể sinh động, khiến cả lớp cười rộ lên khi ông lồng vào những câu chuyện kỳ lạ về quá khứ. Và cuối cùng, tiết Ngoại ngữ sôi nổi cùng cô giáo trẻ vui tính, khiến Suzuki thích thú ra mặt mỗi khi trả lời đúng câu hỏi, dù rằng những gì cô làm chỉ là đang đóng vai Reiko.

Không khí trong lớp thật yên bình. Những chiếc quạt trần quay chậm rãi, cửa sổ mở hé để đón làn gió non. Không có cảnh chen chúc, không có sự ồn ào, chỉ có âm thanh của phấn chạm bảng, của lá cây xào xạc ngoài sân.

Nhưng trong mắt ba người nhà Hikawa, khung cảnh ấy quá quen thuộc. Từng ánh sáng xuyên qua cửa kính, từng âm thanh, từng câu giảng, y hệt như trong vòng lặp đầu tiên.

Từng chi tiết nhỏ bé đều khớp đến mức rợn người, như thể thế giới đang bị tua lại từ đầu, chỉ thay lớp vỏ bề mặt.

``Game-san,'' Suzuki thì thầm, mắt vẫn dán vào bảng đen, ``cháu bắt đầu thấy rõ đường viền của `bối cảnh' rồi. Như thể có ai đang cắt dán từng lớp phông xanh vậy.''

Tiếng đồng hồ trên cổ tay Kuon rung nhẹ, giọng Game chỉ hừm một tiếng, rồi tiếp tục quan sát xung quanh.

Tiếng chuông giải lao vang lên, giai điệu nhẹ nhàng ấy khiến Kuon nhíu mày.

Cậu nhớ rất rõ, nó cũng từng vang lên như vậy, trong một buổi sáng tưởng như bình thường. Và sau giai điệu đó, mọi thứ bắt đầu lặp lại, rồi kết thúc bằng sự sụp đổ.

``\dots Lại khúc dạo đầu,'' Kuon lẩm bẩm, ngón tay siết nhẹ bút chì. ``Lần trước, chỉ sau tiếng chuông này năm phút, cả lớp sẽ đồng loạt quay sang nhìn tôi, và\dots\ ngay bây giờ.''

Cả lớp đồng thanh cất tiếng chào giáo viên rồi quay phắt về phía cậu, mười lăm đôi mắt như được căn chỉnh bằng thước. Kuon nuốt khan.

Kaigou khẽ cười, nhưng mép môi run rẩy: ``Tôi cũng nhớ kịch bản này. Tiếp theo, cô chủ nhiệm sẽ hỏi `Haragen Kyuen-san, em có muốn phát biểu không?' và cậu sẽ đứng dậy, nói câu `Không, thưa cô'.''

Đúng lúc đó, cô giáo gõ thước lên bàn: ``Haragen Kyuen-san, em có muốn chia sẻ với cả lớp không?''

Kuon đứng dậy, giọng khô khốc: ``Không, thưa cô.''

Cậu ngồi xuống, mồ hôi lạnh chảy dọc thái dương. Cô chị gái lẩm bẩm gần như không thành tiếng: ``Chúng ta đang đi đúng hướng. Quả thật, đây đúng là một vòng lặp.''

``Như thể đây là một chương trình được dựng sẵn vậy,'' cậu con trai đáp, mắt nhìn thẳng vào bảng như đang đọc một dòng code chỉ mình thấy. ``Giờ ta phải quan sát xem Shino-san sẽ xử trí như thế nào.''

Cả lớp bắt đầu nhộn nhịp. Một vài nhóm bạn tụ tập bên cửa sổ, số khác rủ nhau xuống căng-tin, còn Shino thì đang nói chuyện cùng với hai người bạn khác: Kana và Kanade.

Nhóm Hikawa nhìn nhau, rồi nhập cuộc, như thể chỉ là những người bạn học bình thường.

``Cậu đã quen với môi trường ở đây chưa, Shino-san?'' Kanade hỏi, giọng nhẹ nhàng. Cô giương đôi mắt lấp lánh sự quan tâm, trong khi đang kiễng người xuống khoanh tay trên bàn Shino.

Kana liền cau mày, chống hông: ``Kanade, bà biết hôm nay mới là ngày thứ hai của cô ấy chứ? Làm gì mà quen được nhanh vậy.''

Kaigou bật cười khẽ, nhanh chóng nhập vao: ``Đúng đấy. Cậu nghĩ người ta có thể hòa nhập ngay à? Ít nhất cũng phải để người ta làm quen dần chứ.''

Shino chỉ biết cười gượng, vài giọt mồ hôi lấm tấm trên trán vì câu hỏi vu vơ đó.

\img{e3-4b}

``Nhắc mới nhớ,'' Kanade khẽ nói tiếp, ``mình nghĩ bọn mình vẫn chưa giới thiệu kỹ với nhau nhỉ.''

Kana gật đầu: ``Chuẩn luôn! Ở trên lớp, tụi mình chỉ nói tên thôi, Shino chắc chưa nhớ ai với ai đâu.''

\begin{wrapfigure}[13]{r}{8cm}
  \includegraphics[height=8cm]{../../../../Image/Book?/kana.png}
\end{wrapfigure}

Hai cô gái quay sang, như thể cùng chia nhau phần giới thiệu bản thân. Trong lúc họ còn đang say sưa, cả ba người Hikawa vừa ngước nhìn họ, vừa nhỏ giọng bàn chuyện riêng, vừa đủ để không ai nghe thấy.

\chrint

``Hmm\dots\ Tokiwa Kana (\ruby{常盤}{ときわ}\ruby{香}{か}\ruby{菜}{な}, \ruby{Tô-ki-oa}{Thường Bàn}\ \ruby{Ca}{Hương}-\ruby{na}{Thái}),'' giọng Game vang lên đầy chậm rãi.

Cô nàng có mái tóc trông như lò xo kia là thành viên thuộc câu lạc bộ phát thanh của trường Aogami. Năng động, hoạt bát và thích giao tiếp, cô luôn tràn đầy năng lượng. Cô và Kanade là bạn từ thuở nhỏ, gắn bó đến mức hầu như lúc nào cũng kề vai sát cánh.

Kana thường xuyên mời Kanade tham gia các chương trình nội bộ của trường với tư cách khách mời. Dù sở hữu tố chất lãnh đạo, cô vẫn có một góc nhỏ mỏng manh trong lòng, điều mà cô chưa từng để lộ ra ngoài.

``\dots Sao giọng cô ấy nghe có vẻ quen thế nhỉ?'' Kuon khẽ nghiêng đầu, mắt vẫn dán vào chiếc loa phát thanh treo tường. ``\textit{Như kiểu\dots\ mình từng nghe cô ấy dẫn chương trình cho trường học ở vòng lặp trước vậy.}''

``Và người kia\dots\ là Amano Kanade (\ruby{雨}{あま}\ruby{乃}{の}\ruby{奏}{かなで}, \ruby{A-ma}{Vũ}-\ruby{nô}{Nãi}\ \ruby{Ca-na-đê}{Tấu}).''

\begin{wrapfigure}[12]{l}{4cm}
  \includegraphics[height=8cm]{../../../../Image/Book?/kanade_human.png}
\end{wrapfigure}

Theo Kana, Kanade là người ít nói nhưng lại sở hữu tấm lòng hiền hậu hiếm có.

Từ thuở nhỏ, hai người đã thân thiết; Kanade luôn ngưỡng mộ sự hoà đồng của Kana và mong rằng mình cũng có thể tự tin đứng trước đám đông như cô bạn.

Điều cô chưa nhận ra là kỹ năng nói trước công chúng của mình đang thăng tiến đáng kể sau mỗi lần được mời làm khách mời trong phòng phát thanh.

\chrint

``Kana\dots\ và Kanade,'' Kuon thì thầm, ngay lập tức cảm nhận được luồng ký ức mơ hồ chạy dọc sống lưng, về hai giọng nói trong bản tin của họ ở vòng lặp trước đó. ``Bảo sao mà tôi thấy quen thế.''

``À mà này,'' giọng của Kana đưa cả nhóm Hikawa về với cuộc trò chuyện, cô nàng tóc tím chỉ sang cô bạn bên cạnh, ``Kanade thực ra là lớp trưởng lớp mình đấy!''

``Đ-Đừng nói linh tinh thế!'' Kanade đỏ mặt, khẽ đập vào vai cô bạn. ``Giới thiệu bản thân bà đi chứ!''

Kana bật cười: ``Chính bà là người tình nguyện làm lớp trưởng để tự tin hơn còn gì.''

Kanade im lặng vài giây, rồi khẽ gật đầu, không phản bác.

Trong khoảnh khắc ấy, Kuon thoáng nhìn Kaigou. Cả hai đều nhận ra điều gì đó.

Một chút Inari trong sự dịu dàng của Kanade. Một chút Hanabi trong sự thẳng thắn của Kana.

Những mảnh ký ức cũ lại chồng lên lớp thực tại mới, khiến hình ảnh trước mặt họ bắt đầu nhiễu nhẹ, méo mó như màn hình TV mất tín hiệu.

Cả hai người họ đồng thời cảm thấy một luồng khí lạnh nhẹ chạy qua cánh tay. Không phải gió, mà là cảm giác quen thuộc đến kỳ lạ. Họ nhìn nhau, không cần nói cũng hiểu: hai cô gái trước mặt chỉ khác biệt ở tên gọi và màu tóc; còn thái độ, nhịp nói, cách họ đứng nghiêng vai về phía nhau\dots\ tất cả đều là bản sao hoàn hảo của Inari và Hanabi. Một thứ ``bình mới rượu cũ'' đang được rót đầy lần nữa, mùi hương cũ thoang thoảng khiến dạ dày họ se lại.

``Các cậu đã biết về bọn mình rồi,'' Kana tiếp, ``giờ đến lượt tụi này tìm hiểu các cậu chứ ha?''

Kaigou tỏ vẻ cười xòa, tiếp tục với màn kịch của mình: ``Ôi dào, bình thường thôi mà. Cứ thoải mái.''

Kanade do dự, nhưng rồi khẽ gật đầu:
``\dots Cậu nói cũng đúng.''

Ngay khi ấy, khung cảnh lại rung lên khẽ khàng.

Trước mắt Kuon và Kaigou, hình ảnh Kanade và Kana chồng lẫn với Inari và Hanabi - hai gương mặt đã từng cười nói như thế, ở một nơi đã bị xóa bỏ.

\begin{figure}[H]
  \centering
  \begin{subfigure}[t]{0.45\textwidth}
    \centering
    \includegraphics[width=8cm]{../../../../Image/Book?/e3-4c1.png}
  \end{subfigure}
  \quad
  \begin{subfigure}[t]{0.45\textwidth}
    \centering
    \includegraphics[width=8cm]{../../../../Image/Book?/e3-4c2.png}
  \end{subfigure}
\end{figure}

``Vậy\dots'' giọng của Hanabi vang lên, dù họ biết đó chỉ là ký ức, ``món ăn yêu thích của các cậu là gì? Môn thể thao yêu thích? Trò chơi điện tử yêu thích? Hay điểm nào khiến người khác thấy thu hút nhất?''

Inari khẽ kéo tay áo cô bạn: ``H-Hanabi-chan\dots\ không cần hỏi đến thế đâu\dots''

``Sao vậy? Chẳng phải chúng ta nên hiểu rõ về họ hơn sao?''

``Thì, đúng là vậy\dots nhưng đâu cần hỏi điểm thu hút của họ\dots''

Hình ảnh nhòe dần, rồi trở lại thực tại. Kana đang cười, Kanade vẫn hơi bối rối.

``Sao thế, Shino-san?'' cô nàng tóc tím nghiêng đầu, nhìn Shino đang trơ mặt ra.

Cô ấy liền khẽ lắc đầu, mỉm cười, như thể cô đã đoán định được trước: ``Không\dots\ không có gì đâu. Chỉ là\dots\ mình đang tưởng tượng cái gì đó thôi.''

Không ai trong nhóm Hikawa đáp lại.

Nhưng trong ánh mắt họ - ba người đã từng trải qua những kết thúc tan vỡ - câu nói ấy như một lời cảnh báo.

Bởi họ biết rõ hơn ai hết: chính thế giới này mới là thứ được tưởng tượng ra.

Ngay khi mọi người còn chưa hết sững sờ vì đoạn nhiễu vừa tan, ba giọng nói khác lại vang lên, đan xen nhau giữa lớp học đầy ánh nắng.

``Cô bạn học sinh mới chuyển kia tưởng thế mà cũng thu hút cơ đấy,'' giọng một cô gái cất lên.

``Chẳng phải điều đó là hiển nhiên sao\~{}'' một giọng khác cười khúc khích, giọng hơi ngân ra.

``Nhưng nếu chúng ta cứ bâu vào cô ấy như vậy, không sợ cậu ấy bị ngợp sao?'' cô thứ ba nhẹ giọng xen vào.

Suzuki ngoảnh lại về hướng phát ra tiếng nói, và tim cô chợt thắt lại. Đứng trước mặt cô là ba gương mặt quá đỗi quen thuộc: Saya, Kaoru và Miku.

``Có chuyện gì mà mọi người bàn tán xôn xao thế?'' Kaoru nghiêng đầu hỏi.

``Cho bọn này tham gia với\~{}'' Saya nhoẻn cười.

Thế nhưng khi lớp học lại chập chờn, hình ảnh trước mắt cô mờ đi rồi tái hiện. Saya, Kaoru, Miku biến mất, thay vào đó là ba người khác: Himekawa Mitsuki, Shiratori Hotaru, và Kaizuka Yuki.

\begin{figure}[H]
  \centering
  \begin{subfigure}[t]{0.45\textwidth}
    \centering
    \includegraphics[width=8cm]{../../../../Image/Book?/e3-4d1.png}
  \end{subfigure}
  \quad
  \begin{subfigure}[t]{0.45\textwidth}
    \centering
    \includegraphics[width=8cm]{../../../../Image/Book?/e3-4d2.png}
  \end{subfigure}
\end{figure}

Cô không thốt nên lời. Cử chỉ, ánh mắt, cả giọng nói của họ\dots\ đều mang một cảm giác trùng lặp rợn người: Mitsuki có lối ngân giọng như Saya, Hotaru thì thẳng tính chẳng khác gì Kaoru, còn Yuki lại toát ra nét dịu dàng mà Miku từng có.

\chrint

\begin{wrapfigure}[14]{r}{6cm}
  \includegraphics[height=8cm]{../../../../Image/Book?/hotaru_human.png}
\end{wrapfigure}

Shiratori Hotaru (\ruby{白}{しら}\ruby{鳥}{とり}\ruby{火}{ほ}\ruby{樽}{たる}, \ruby{Si-ra}{Bạch}-\ruby{tô-ri}{Điểu}\ \ruby{Hô}{Hỏa}-\ruby{ta-rư}{Tôn}), cái tên vừa vang lên đã khiến cả lớp im phăng phắc để chờ đợi phán quyết thẳng thừng. Thẳng thắn, đáng tin, cô là điểm tựa mỗi khi ai đó cần lời khuyên không đường vòng; cũng là trung tâm chú ý khi tiếng cười giòn tan của cô bất chợt cất lên giữa giờ giải lao.

Tuy vậy, ẩn sau vẻ cứng cỏi ấy lại là căn phòng ngập thú nhồi bông: gấu trúc nằm gọn trong lòng, mèo máy xếp hàng trên kệ, cá heo bé xíu được ôm riêng vào đêm, những thứ ``dễ thương'' khiến cô đỏ tai mỗi lần bắt gặp ánh mắt soi mói. Khi ai đó buông một lời khen ``Hotaru dễ thương quá'', cô lại lúng túng, nói ấp úng ``Đâu\dots\ đâu có'', rồi giấu mặt vào khăn quàng, để lại phía sau hình ảnh trái ngược hoàn toàn với cô bé ``thép'' thường ngày.

Kana từng mời cô làm khách mời đài truyền thanh; buổi phát sóng hôm ấy gần như chuyển thành ``chương trình của Hotaru'' khi cô nói liên tục từ đầu đến cuối, từ kỷ niệm tiểu học đến cách chọn bút chì màu, khiến Kana phải giơ biển ``NGỪNG'' ba lần. Sau giờ phát sóng, cô bị mắng te tua, cúi đầu xin lỗi rồi chạy về lớp, mặt đỏ như gấc.

\begin{wrapfigure}[15]{l}{6cm}
  \includegraphics[height=8cm]{../../../../Image/Book?/yuki.png}
\end{wrapfigure}

Kaizuka Yuki (\ruby{械}{かい}\ruby{塚}{ずき}\ruby{幸}{ゆき}, \ruby{Cai}{Giới}-\ruby{dư-ki}{Trũng}\ \ruby{Giu-ki}{Hạnh}) là kiểu học sinh ``thích ở bên mọi người hơn cả việc ăn ba bữa một ngày''; với cô, những giờ học tập ở trường còn vui hơn cả ngày nghỉ, thậm chí còn coi``trường học là ngôi nhà thứ hai'' theo đúng nghĩa.

Giọng hát trong trẻo của cô vang lên hầu như mỗi buổi trưa, khi Yuki rủ cả lớp đi karaoke ngay trong phòng học trống. Ban đầu, bạn bè còn e ngại, dần dần từ chối, nhưng cô không nản; thay vào đó, cô biến giờ nghỉ thành buổi diễn thường nhật, ép mọi người cầm micro cùng hát. Trong khi các bạn còn bối rối, Yuki nở nụ cười mãn nguyện, tuyên bố rằng khả năng ca hát của cả lớp đang tiến bộ từng ngày, và còn tuyên bố rằng ``chỉ cần thêm vài bài nữa, họ sẽ thành dàn hợp xướng chuyên nghiệp của trường Aogami đó\~{}''.

\begin{wrapfigure}[14]{r}{7cm}
  \includegraphics[height=8cm]{../../../../Image/Book?/mitsuki.png}
\end{wrapfigure}

Himekawa Mitsuki (\ruby{姫}{ひめ}\ruby{川}{かわ}\ruby{美}{み}\ruby{月}{つき}, \ruby{Hi-mê}{Cơ}-\ruby{ca-oa}{Xuyên}\ \ruby{Mi}{Mỹ}-\ruby{xư-ki}{Nguyệt}), một cái tên nghe nhẹ như làn gió đầu thu, nhưng ẩn sau vẻ ngoài mềm mại ấy lại là một trái tim rực lửa không kém ai. Cô gái ấy sở hữu đôi mắt nâu ấm, luôn hơi nheo khi cười, khiến người đối diện lầm tưởng nàng thuộc tuýp ``dâu tằm'' dịu dàng, dễ gãy.

Thực ra, Mitsuki chỉ đang\dots\ giữ năng lượng. Bước chân cô vội vã trên phố, tiếng guốc gõ đều đặn như beat nhạc dance, mới là bản năng thật: thích nơi đông người, thích tiếng rao của vendor, thích mùi bánh crepe nóng hổi trước ga Aogami, nơi cô thường ``tình cờ'' xuất hiện vào mỗi chiều thứ Năm, tay xách bao ni-lông đủ màu, miệng cười tươi như vừa thắng một ván trò chơi.

Những ai từng bắt gặp Mitsuki và Kanade cùng nhau sẽ nghĩ họ là đôi bài trùng hợp: hai nụ cười sáng, hai mái tóc bay, hai cái vẫy tay cùng nhịp độ. Nhưng khi bạn thân rời đi, để lại Mitsuki một mình trước cửa hiệu phụ kiện, nàng không hề rơi vào khoảng lặng buồn bã; thay vào đó, cô lập tức bật chế độ ``thám hiểm'', lướt qua từng quầy hàng, thử băng đô, soi gương, gật gù tự kỷ, hoàn toàn tự chủ và vui vẻ. Người qua đường thấy thế, lắc đầu cười: ``Himekawa-san đúng là con ngựa bất kham trong lồng nhung.''

\chrint

Suzuki khẽ kéo tay áo Kaigou, lí nhí: ``Onee-sama\dots\ em thấy ba người kia\dots\ sao mà quen quá.''

Kaigou chau mày, mắt vẫn dõi theo bộ ba đang tíu tít trước cửa hàng. ``Ừ\dots\ như từng gặp ở\dots''

``\dots trên chuyến tàu vòng lặp trước,'' Kuon thì thầm tiếp lời, khớp hàm cứng lại. ``Cuốn kỷ yếu.''

Cả ba cùng nhớ ra: tàu điện ngầm chật cứng, tiếng loa rè rè, và cuốn kỷ yếu vất vưởng trên ghế, và ngoài những chân dung bị gạch xóa, còn có ảnh ba khuôn mặt này, cùng với bản thân họ và Shino.

``Mà này,'' Mitsuki quay sang Hotaru, ``cậu có nghe chuyện gì bên lớp kế bên chưa?''

Yuki khẽ trầm giọng: ``Ừm\dots\ gần đây càng ngày càng có nhiều bạn học vắng mặt. Lớp bên cạnh yên tĩnh hẳn luôn ấy\dots''

``Với cả,'' Mitsuki tiếp lời, ``chẳng một ai liên lạc được với họ cả.''

Hotaru thở dài: ``Chắc là cúm hè thôi mà?''

Ba người Hikawa nhìn nhau. Shino đột nhiên im lặng, ánh mắt dán chặt xuống bàn. ``\hl{Cô ấy đang cảm thấy quen thuộc,}'' Game thì thầm qua giọng kim loại lạnh lẽo. ``\hl{Còn các cậu, hẳn cũng biết rõ rồi.}''

Cả ba người gật đầu, hiểu được ý ông ta đang nói gì.

``\hl{Bọn tôi cũng đã từng nghe việc Inari, Touka và Saki nói thế từ vòng lặp đầu tiên.}'' họ cùng đồng thanh. ``\hl{Những `vụ mất tích' này đều bắt nguồn từ ga Aogami.}''

Kuon cảm thấy một luồng lạnh chạy dọc sống lưng, như thể có ai đang rút dần không khí khỏi phổi cậu. Cậu nhìn xuống bàn tay mình: năm đầu ngón rung nhè nhẹ, không phải vì sợ, mà vì cái cảm giác ``đã biết trước'' đang bóp nghẹt.

Kaigou khoanh tay trước ngực, móng tay chìa vào da thịt. Cô cảm thấy một mùi sắt tan thoang thoảng, có lẽ do vết máu cũ thoang thoảng ở ddaua đó. Trong đầu, những mảnh ghép của các vòng lặp trước xoay tròn như chiếc đĩa bị trầy xước: Inari mỉm cười, Touka gật đầu, Saki vẫy tay, rồi tất cả đều bỏ mạng ngay trước mắt và biến thành những cái xác vô hồn chỉ lặp đi lặp lại.

Suzuki thì lặng lẽ siết chặt nắm tay lại. Cô cảm thấy thời gian như đặc quánh quanh mình: mỗi giây trôi qua là một bước chân xuống cầu thang vô hình, không có điểm kết thúc. Cô thầm đếm nhịp thở: từng giây hít vào, từng giây giữ lại, từng giây thở ra, không để thực tại nuốt chửng mình trước khi mình kịp nhìn thấy nó.

Cả ba không cần trao đổi ánh mắt; họ đã đủ hiểu nhau sau bao vòng lặp. Họ đang cùng chạy trong một mê cung mà bản đồ được vẽ bằng những sự kiện rối rắm mỗi lần lặp lại, đường đi thay đổi một chút, nhưng điểm kết thúc vẫn là cái nỗi sợ không tên. Và giờ đây, khi những cái tên quen thuộc lại vang lên, họ biết: vòng lặp mới đã chính thức quay, và lần này, họ không được phép chỉ là khán giả.

``Vậy\dots\ liệu họ có biết về việc đó không ạ?'' cô em út liếc nhìn họ, đôi mắt hiện rõ vẻ e sợ.

``\hl{Tôi e là không,} giọng Pháp trầm đáp lại, tỏ vẻ ngán ngẩm. ``\hl{họ đều đã bị lập trình sẵn để nói ra những câu đó.}''

Kuon ngẩng lên. Mọi thứ xung quanh họ đang thay đổi. Kính cửa sổ nứt vỡ, vài chỗ còn được chắp vá bằng tấm gỗ xám xịt. Những bạn học xung quanh không còn mặc đồng phục mùa hè nữa, thay vào đó là áo blazer đen, áo len cổ cao, và cả hơi lạnh lạ lùng lan trong không khí.

\begin{figure}[H]
  \centering
  \begin{subfigure}[t]{0.45\textwidth}
    \centering
    \includegraphics[width=8cm]{../../../../Image/Book?/e3-4e1.png}
  \end{subfigure}
  \quad
  \begin{subfigure}[t]{0.45\textwidth}
    \centering
    \includegraphics[width=8cm]{../../../../Image/Book?/e3-4e2.png}
  \end{subfigure}
\end{figure}

``\textit{Thời gian trong thế giới này đã bị đảo lộn hoàn toàn\dots}'' Kaigou thầm nghĩ, ``\textit{như thể nơi này đang dần bộc lộ bản chất thật của mình.}''

Cả lớp giờ chỉ chú ý vào một người duy nhất: Misora Shino, người vừa mới chuyển tới.

Hotaru nghiêng người, hỏi với nụ cười: ``Vậy thì, người bạn mới chuyển à. Cậu đã quen với môi trường này rồi chứ?''

``Ể!? M-Mình quen rồi!'' Shino giật mình đáp, gượng cười.

``May quá đi\~{}'' Yuki mỉm cười hiền hậu.

``Nếu có chuyện gì cần giúp thì cứ nói với bọn mình nha\~{}'' Mitsuki nói thêm, rồi liếc sang Kanade, nở nụ cười tươi: ``Đúng không nào, Inari-san?''

Khoảnh khắc cái tên Inari vang lên, Kuon, Kaigou và Suzuki cùng sững người. Tất cả nghi ngờ giờ đã thành chắc chắn: thế giới ``hạnh phúc'' này chỉ là một vòng lặp, một chương trình được dựng lại để trói buộc ký ức của họ.

Suzuki khẽ nói, giọng run: ``Vấn đề là\dots\ liệu Shino-san có nhận ra được không?''

``Tôi nghĩ,'' Game đáp khẽ, ``chúng ta chỉ còn có thể chờ thôi.''

Ở phía xa, Shino tươi cười giữa đám bạn mới, giọng cô hồn nhiên vang lên: ``Mình muốn tìm hiểu rõ hơn về thị trấn Aogami này! Mình cũng muốn biết thêm về các cậu nữa!''

Trước mắt ba người Hikawa, cô ấy được năm bạn học vây quanh, trò chuyện ríu rít, một khung cảnh ngập nắng nhưng lại phảng phất thứ gì đó giả tạo, như thể tất cả niềm vui ấy đều được lập trình từ trước.

\begin{dialogue}
  \speak{Kana} Shino-san à, bọn mình muốn cậu làm khách mời cho buổi thông báo trưa nay nhé!
  \speak{Kanade} Mình biết vài chỗ hay lắm. Nếu cậu muốn, mình dẫn đi liền!
  \speak{Hotaru} Ơ\dots gần đây có một cửa hàng bán thú nhồi bông dễ thương cực luôn đó\dots
  \speak{Mitsuki} Quanh nhà ga có nhiều thứ thú vị lắm nha\~{} Đi cùng bọn mình đi!
  \speak{Yuki} Thôi nào, từng người nói một thôi! Các cậu làm Shino-san rối đấy! \dots À mà này, cậu có thích ca hát không?
\end{dialogue}

\img{e3-4f}

Trong khi Shino cười rạng rỡ giữa vòng vây bạn bè, ba người còn lại chỉ biết đứng nhìn, nụ cười của cô phản chiếu trên mặt kính nứt vỡ, méo mó và mong manh như chính thực tại này.

``Ôi\dots\ mình thấy hạnh phúc quá\dots'' cô thì thầm, đôi mắt khép hờ trong khoảnh khắc, như thể cả thế giới này chỉ còn lại tiếng cười của cô và những bông hoa ảo đang nở quanh.

Kuon siết chặt quai cặp, môi mím đến trắng bệch; Kaigou đứng thẳng như tượng, ánh mắt lạnh toát; Suzuki chúc mũi, bàn tay vô thức nắm lấy góc áo chị.

Ba bóng dáng cứng đờ giữa dòng người đang xoay tròn, như ba mảnh đá chìm trong ly nước đường ngọt sắc, không thể nổi, không thể tan, chỉ có thể im lặng mà chứng kiến niềm vui giả tạo ấy lan tỏa.

Tiếng chuông ngân vang, kéo dài, vang vọng như vọng lại từ một nơi xa xôi nào đó: chậm hơn, trầm hơn, và có phần nghèn nghẹt như phát ra từ một chiếc loa cũ kỹ. Cả lớp nhốn nháo đứng dậy, tiếng ghế kéo lạch cạch hòa vào những lời chào tạm biệt ríu rít. Ánh nắng buổi chiều hắt qua khung cửa gỗ, vàng vọt và dày đặc bụi, phủ lên tất cả một sắc ấm đến giả tạo.

Thời gian trôi chậm lại, hay đúng hơn, bị kẹt lại. Ánh sáng không thay đổi. Âm thanh ngoài cửa sổ như bị nuốt chửng. Chiếc đồng hồ treo tường chỉ nhích kim một lần, rồi dừng lại ở vạch giữa hai con số.

Shino vẫn ngồi đó, bàn tay khẽ siết lấy vạt áo. Một phần trong cô muốn tin rằng mọi thứ hôm nay thật đẹp, những nụ cười, những lời mời rủ, cái cách ai cũng đối xử với cô bằng sự ấm áp. Nhưng\dots\ càng cố gắng tin, lòng cô càng thấy lạnh. Có điều gì đó lệch nhịp, như thể cô đang diễn lại một vở kịch mà ai đó đã viết sẵn từ trước.

Một làn gió nhẹ thổi qua khung cửa nứt, làm lay động tờ giấy dán trên bảng: dòng chữ ``Hãy cùng nhau tận hưởng khoảng thời gian học đường này nhé!'' khẽ rung lên, trước khi rơi xuống sàn.

Kuon, Kaigou và Suzuki vẫn chưa rời đi. Họ đứng phía cuối lớp, im lặng nhìn Shino.

``\vie{Anh nghĩ,}'' cậu con trai khẽ nói, “dù cho có bao nhiêu lớp hạnh phúc được dựng lên\dots\ cảm giác thật vẫn không thể giấu được mãi.''

Cô em út im lặng, mắt vẫn dõi theo Shino, cô gái đang ngồi đó, giữa căn phòng vàng nhạt, giữa thế giới được lập trình để cười.

Cô chị lớn gật nhẹ: ``\vie{Nếu cô ấy còn cảm thấy điều gì đó lạ lẫm\dots\ nghĩa là vẫn còn lối thoát.}''

Tiếng cửa lớp khẽ kêu cạch khi họ ra về, để lại Shino một mình trong căn phòng đã bắt đầu tối dần. Cô nhìn quanh, và lần đầu tiên, nhận ra rằng chiếc bóng của mình trên nền gạch\dots\ không di chuyển, dù ánh sáng ngoài cửa sổ đã thay đổi.

Một nụ cười thoáng hiện nơi khóe môi. Một nụ cười mảnh mai, run rẩy, và gần như tuyệt vọng.

``Hạnh phúc thật lạ\dots'' Cô thì thầm. ``Nó ấm áp\dots\ nhưng lại khiến mình lạnh người.”

Ngoài kia, tiếng chuông báo hết giờ lại vang lên một lần nữa, y hệt như tiếng chuông lúc trước, chẳng nhanh hơn, chẳng chậm hơn. Và cứ thế, vòng lặp trong thế giới hạnh phúc vẫn cứ tiếp tục\dots


\end{document}